
Maximum power point tracking (MPPT) changes the operating point of a photovoltaic source to increase extracted power. A data-driven controller can estimate a reference directly or learn a current--voltage characteristic from which a reference is computed. These tasks are related but not interchangeable. A smaller current prediction error does not imply that the location of the predicted power maximum becomes more accurate.

Neural MPPT has already been evaluated on embedded hardware. Abouzeid et al. compared a neural method with perturb and observe (P\&O) and incremental conductance (INC), including experimental implementation on a low-cost microcontroller \cite{abouzeid}. Incremental learning has also been investigated. Amine et al. described a Raspberry Pi-based meteorological acquisition workflow that expands the dataset and supports updates of a neural estimator \cite{amine}. Jamshed et al. reported a lightweight neural controller implemented on an STM32 and evaluated through controller hardware-in-the-loop testing \cite{jamshed}. These studies establish that embedded inference, model updating and controller validation each have prior examples. They also motivate an explicit distinction between computation on a physical controller and execution within an electronic simulator.

The present work builds on the sensing-aware simulation-based digital twin framework developed by Suárez Suárez et al. \cite{suarez}. That publication supplies the methodological origin of the digital photovoltaic plant and sensing chain. The present study examines an additional controller, its bounded online identification process and the difference between predictive improvement and energy tracking. The plant equations are retained. The experimental control interface is documented separately from the original project.

The research question is whether updating a small current predictor during operation produces a measurable energy benefit compared with the same network with frozen weights. A paired comparison is needed because a functioning controller and a changing parameter vector do not establish such a benefit. Conventional algorithms provide context for the tested implementation. The frozen-network ablation isolates the effect of online updates more directly.

